\section{Continuum barriers and lattice traversal counts}
\label{sec:continuum-transfer}

Traversal counts separate the large-scale accumulation of excursions from
occupation inside the last annulus.  We first record the continuum barrier
estimates in the form needed below.  We then compare rare upper events by
multiplying one-step lattice probabilities.  For the lower event, a single
strong coupling transfers the existence of a suitable Brownian center.
The two comparisons have different purposes: the first preserves probabilities
on their exponential scale, while the second preserves a positive probability
of existence.
Together they give the three lattice inputs used later: an upper bound on
barrier crossings, a joint bound for a barrier-respecting traversal endpoint,
and a positive probability of finding a lower root.  The first two enter the
localized upper bound; the third enters the lower bound for the disk maximum.

\subsection{Spherical traversal estimates}

We use Rosen's spherical traversal construction \cite{Rosen2022}.
The radius map $h(r)$ below is distinct from the scalar terminal depth
$h=\log M$ used later in the lattice comparison.
Its upper and lower arguments use separate fixed choices of the initial
radius.  We keep the customary notation $r_0$ in each argument and state
the change explicitly; the two choices are never used in the same pathwise
comparison.

Let $X$ be spherical Brownian motion started at $v$, stopped at its
first hit $\tau$ of $\partial B_d(v,r_*)$.  When defining source
stopping times we extend it by an independent Brownian continuation;
the safety event below will imply that every portion actually used lies
before $\tau$.  Put
\begin{equation}
h(r)=2\arctan(r/2),\qquad r_j=r_0e^{-j},\qquad h_j=h(r_j),\label{eq:ct-external-23}
\end{equation}
For the upper estimates, choose the fixed radii as in
\cite[Section~2, (2.2)]{Rosen2022}, in particular $2r_*\le h_0$.  Let $F_j$ be a
maximal $d_0h_j$-separated set, augmented if necessary to a
$d_0h_j$-cover; thus $|F_j|\asymp e^{2j}$.  For $y\in F_L$,
$T_{y,j}$ is the number, before $\tau$, of completed excursions from
$\partial B_d(y,h_{j-1})$ to $\partial B_d(y,h_j)$, with the initial
and final incomplete pieces omitted.  Its source shell is
\begin{equation}
k_y=\inf\{k:y\notin B_d(v,h_k)\}.\label{eq:ct-external-24}
\end{equation}
Pairs have separation shell $k$ when
$2h_k<d(y,y')\le2h_{k-1}$.

Set
\[
\rho_L=2-\frac{\log L}{L},\qquad l_L=l\wedge(L-l),
\]
\begin{equation}
\alpha_{z,+}(l)=\rho_Ll+z+l_L^{1/4},\qquad
 \alpha_{z,-}(l)=\rho_Ll+z-l_L^{1/4}.\label{eq:ct-external-25}
\end{equation}
Fix an admissible $z\ge z_0$ in Rosen's upper argument.  The upper
first-crossing event $U(y,l,z)$ means
\[
\sqrt{2T_{y,j}}\le\alpha_{z,+}(j)
       \quad(k_y<j<l),\qquad
 \sqrt{2T_{y,l}}>\alpha_{z,+}(l).
\]
Rosen's representative construction is deterministic.  With fixed
tie-breaking, every fine center is sent to its closest point of the
predetermined coarser net at scale
$h_{l+\frac14 l_L^{1/4}}$; equations (2.26)--(2.33) give the buffered
ball inclusions and domination of its traversal count.  Proposition 2.5,
specifically equations (2.34)--(2.35), then proves the complete estimates
\begin{equation}
\#\{\hbox{representatives with source }k\}
 \le C\exp\{2(l-k)+\tfrac12l_L^{1/4}\},\label{eq:ct-external-26}
\end{equation}
\begin{equation}
P(U(y,l,z))\le
 C\exp\{-2l-l_L^{1/4}\}e^{-2z}.\label{eq:ct-external-27}
\end{equation}
They hold pointwise for $1\le k\le l\le L-(4\log L)^4$.
Only the choice $F_L\cap B_d(v,h_{\log L})^c$, and not the proof of
\eqref{eq:ct-external-26}--\eqref{eq:ct-external-27}, imposed $k\le\log L$.  Hence, for any
$s\le C_0\log L$, summing just the representatives whose source shell
is at least $s$ gives
\begin{equation}
P\big(\exists y:k_y\ge s,\ \exists l\le L-(4\log L)^4:
              U(y,l,z)\big)
       \le C e^{-2s+o(\log L)}e^{-2z}.\label{eq:ct-external-28}
\end{equation}

\subsubsection{The endpoint estimate for complete centered excursions}

The endpoint calculation must be applied to a complete centered excursion
forest.  To specify that law, start spherical Brownian motion at $v$ with
$h_{k_y}\le d(v,y)<h_{k_y-1}$ and continue it until its first exit from
$B_d(y,h_0)$.  Let $T^{\mathrm c}_{y,j}$ count its completed traversals
at level $j$, and let ${\cal G}^{\mathrm c}_{y,L,z}$ require the upper
barrier in \eqref{eq:ct-external-25} from $m=k_y+1$ through $L-1$.
Put $Y_y^{\mathrm c}=\sqrt{2T^{\mathrm c}_{y,L}}$ and
$B_L=\rho_LL+z$.  Uniformly for
$k_y\le C_0\log L$, every fixed $\epsilon_0>0$,
$\epsilon_0L\le q\le B_L+1$, and fixed $z$, the source law in
\cite[(2.8)]{Rosen2022} and its Appendix Theorem~9.1(a) give
\begin{equation}
P\big({\cal G}^{\mathrm c}_{y,L,z},Y_y^{\mathrm c}\in[q,q+1]\big)
 \le C(k_y+1)^C(2+B_L-q)L^{-2}e^{-q^2/(2L)}.\label{eq:ct-external-29}
\end{equation}
This assertion concerns the centered stopping rule just defined.  We do not
identify the descendants of a population stopped at the off-center boundary
$\partial B_d(v,r_*)$ with an independent Galton--Watson forest.
The later lattice endpoint estimate will be proved directly for its centered
continuation.

The following version of \cite[Appendix, Theorem~9.1(a)]{Rosen2022}
is sufficient: if $T$ is critical Galton--Watson with
$P(\xi=m)=2^{-m-1}$, $T_0=x^2/2$, and
\[
\begin{gathered}
\delta>0,\quad C\ge0,\quad\eta>1,\quad\epsilon\in(0,1/2),\\
2\le x,y\le\eta L,\quad x^2/2\in\mathbb N,\\
0\le x\le a,\quad0\le y\le b,
\end{gathered}
\]
\begin{equation}
\begin{split}
P\big(&\sqrt{2T_j}\le a+(b-a)j/L
             +C(j\wedge(L-j))^{1/2-\epsilon},\ 1\le j<L;\\
&\sqrt{2T_L}\in[y,y+\delta]\big)
\le C'\frac{(1+a-x)(1+b-y)}L\sqrt{\frac{x}{yL}}
e^{-(x-y)^2/(2L)}.\label{eq:ct-external-30}
\end{split}
\end{equation}
Write $m=k_y+1$ and $N=L-m$.  Conditional on $T^{\mathrm c}_{y,m}=t$,
the $t$ complete source excursions give the critical geometric
Galton--Watson chain in \eqref{eq:ct-external-30}, started from $t$.
Indeed each excursion from level $m$ to level $m-1$ is explored completely
before the centered path can exit at level $0$; its neighboring logarithmic
side choices have probability $1/2$ independently of the entrance angle.
The logarithmic
gambler--ruin calculation (2.8) gives
\[
P(T^{\mathrm c}_{y,m}\ge t)\le Cm e^{-t/m},
\]
and hence, after grouping the integer values of $t$ for which
$x=\sqrt{2t}$ lies in a unit band,
\begin{equation}
P(\sqrt{2T^{\mathrm c}_{y,m}}\in[x,x+1])
       \le Cm^C e^{-x^2/(2m)}.\label{eq:ct-external-30a}
\end{equation}
On ${\cal G}^{\mathrm c}_{y,L,z}$, $x\le a:=\alpha_{z,+}(m)=O(m)$.
Apply \eqref{eq:ct-external-30} over the remaining $N$ generations with this $a$ and
with $b=B_L+1$.  The needed deterministic barrier inclusion follows
from
\begin{equation}
((m+j)\wedge(L-m-j))^{1/4}
 \le (1-j/N)m^{1/4}+C(j\wedge(N-j))^{1/4};\label{eq:ct-external-30b}
\end{equation}
for $j\le N/2$ use
$(m+j)^{1/4}\le m^{1/4}+j^{1/4}$ and
$m^{1/4}j/N\le j^{1/4}$, while for $j>N/2$ use the
right-end distance directly.  Thus every parameter in \eqref{eq:ct-external-30} has been
specified, including $x\le a$, $q\le b$, and
$x,q\le\eta N$.

The sole positive start not covered by the hypothesis $x\ge2$ is
$t=1$.  Couple a process from two ancestors as the sum of two
independent one-ancestor processes.  On the event, of probability
$1/2$, that the second ancestor has no child in the first generation,
the total process agrees from generation $1$ onward with the first
process.  The barrier in \eqref{eq:ct-external-30} is imposed precisely at generations
$1,\ldots,N-1$, so the $t=1$ contribution is at most twice the same
event for initial population two, to which \eqref{eq:ct-external-30} applies.  The case
$t=0$ cannot have the positive endpoint $q\ge\epsilon_0L$.

Multiplication of \eqref{eq:ct-external-30a} with \eqref{eq:ct-external-30}, the identity
\begin{equation}
\frac{x^2}{k+1}+\frac{(q-x)^2}{L-k-1}
 =\frac{q^2}{L}+
   \frac{L(x-(k+1)q/L)^2}{(k+1)(L-k-1)},\label{eq:ct-external-31}
\end{equation}
proves \eqref{eq:ct-external-29}.  Indeed $q\ge\epsilon_0L$, $N\asymp L$,
$x\le Cm$, and so the square-root factor in \eqref{eq:ct-external-30}, together with its
explicit $N^{-1}$, is at most $Cm^{1/2}L^{-2}$; all remaining source
and summation factors cost only $m^C$.

No estimate of the form \eqref{eq:ct-external-29} is asserted at $q=0$: extinction there
has order-one probability.  In the lattice upper bound below,
$q<\epsilon_0L$ is instead disposed of by the saturated terminal
Chernoff factor.  We shall also use the following truncated version.
If $n=J-h$, $m\le C_0\log J$, and
\[
B_l=\rho_Jl+z+(l\wedge(J-l))^{1/4},
\]
then \eqref{eq:ct-external-29}, with $L$ replaced by $n$ and $B_L$ by $B_n$, holds
for $\epsilon_0n\le q\le B_n+1$.  To verify the only new point, put
$\phi(l)=(l\wedge(J-l))^{1/4}$.  Its two Hölder bounds from $m$
and $n$ give
\begin{equation}
\phi(m+j)\le(1-j/N)\phi(m)+(j/N)\phi(n)
                 +C(j\wedge(N-j))^{1/4}.\label{eq:ct-external-31a}
\end{equation}
The affine part $\rho_Jl+z$ interpolates exactly, so \eqref{eq:ct-external-30} applies with
the arbitrary endpoints $a=B_m,b=B_n+1$.  In particular, no
$\rho_n$-barrier theorem is being assumed.

\subsubsection{A fixed number of source excursions}

For the lower event, now choose $r_0$ sufficiently small that
$4h(r_{-1})<r_*$, as in \cite[Section~3]{Rosen2022}.  The preceding
upper estimates retain their separate choice of radii.  Fix the source
parameter $x_*>0$ from that section so that
\begin{equation}
m_0:=x_*^2\in\mathbb N.\label{eq:ct-external-32a}
\end{equation}
(This $x_*$ is not the variable called $x$ in Appendix Theorem 9.1;
the latter equals $\sqrt{2m_0}$ when that theorem is started with
$m_0$ particles.)  Put
\[
F_L^0=F_L\cap\{h_1/40\le d(v,y)\le h_1/20\}.
\]
This fixed annulus, instead of Rosen's full ball, retains a fixed
positive area fraction; the one- and two-point sums in Rosen's proof are
therefore unchanged up to constants, and every source index is $O(1)$.

We spell out the superscript in Rosen (3.3), because it is important for
the stopping-time comparison below.  Write
$\Gamma_{y,j}=\partial B_d(y,h_j)$ and set
\[
A_{y,1}=H_{\Gamma_{y,1}},\qquad
 B_{y,i}=\inf\{t>A_{y,i}:X_t\in\Gamma_{y,0}\},
\]
\begin{equation}
A_{y,i+1}=\inf\{t>B_{y,i}:X_t\in\Gamma_{y,1}\},
                    \qquad 1\le i<m_0.\label{eq:ct-external-32b}
\end{equation}
The intervals $[A_{y,i},B_{y,i}]$, and only these intervals, are the
first $m_0$ source excursions; the intervening intervals
$[B_{y,i},A_{y,i+1}]$ are connectors and are not counted.  Set
$T_{y,1}^{1,m_0}=m_0$.  For $l\ge2$, define inside each source
interval
\[
C_{i,l,1}=\inf\{t\ge A_{y,i}:X_t\in\Gamma_{y,l-1}\},
\]
\[
D_{i,l,q}=\inf\{t>C_{i,l,q}:X_t\in\Gamma_{y,l}\},\qquad
 C_{i,l,q+1}=\inf\{t>D_{i,l,q}:X_t\in\Gamma_{y,l-1}\},
\]
and put
\begin{equation}
T_{y,l}^{1,m_0}=\sum_{i=1}^{m_0}
             \#\{q:D_{i,l,q}<B_{y,i}\}.\label{eq:ct-external-32c}
\end{equation}
Thus Rosen's abbreviation is exactly
$T_{y,l}^{\,1}=T_{y,l}^{1,x_*^2}$, not a count in one source
excursion.  The strong Markov property and logarithmic gambler's ruin
make $(T_{y,1+j}^{1,m_0})_{j\ge0}$ a critical Galton--Watson chain
with $m_0$ ancestors and offspring law
$P(\xi=r)=2^{-r-1}$, $r\ge0$.

Define $\widehat I_{y,z}$ by
\begin{equation}
\sqrt{2T_{y,l}^{1,m_0}}\le\alpha_{z,-}(l),\quad1\le l<L,
 \qquad \sqrt{2T_{y,L}^{1,m_0}}\ge\rho_LL+z.\label{eq:ct-external-32}
\end{equation}
For the outward excursions at scale $k$, let $\theta_{k,i}$ be their
angular increments and let $\nu_k$ be the exit-angle law from radius
$r_k$ to $r_{k-1}$.  If
\[
N_{k,a}=\left\lfloor\frac{(\rho_Lk+z-a+1)^2}{2}\right\rfloor,
\]
 set $N_k=N_{k,a}$ on
$\sqrt{2T_{y,k}^{1,m_0}}\in[\rho_Lk+z-a,\rho_Lk+z-a+1]$, and
\begin{equation}
W_{y,k}(N_k)=\left\{
 d_{W^1}\left(N_k^{-1}\sum_{i\le N_k}\delta_{\theta_{k,i}},
                 \nu_k\right)
 \le\frac{c_0\log(L-k)}{2\sqrt{N_k}}\right\}.\label{eq:ct-external-33}
\end{equation}
With $L^+=L-(500\log L)^4$ and Rosen's fixed $d_*$, the event is
\begin{equation}
I_{y,z}=\widehat I_{y,z}
          \cap\bigcap_{k=L^+}^{L-d_*}W_{y,k}(N_k).\label{eq:ct-external-34}
\end{equation}
Rosen Lemmas 3.3--3.4 state completely that, uniformly for
$0\le z\le\log L$,
\begin{equation}
P(I_{y,z})\asymp(1+z)e^{-2L}e^{-2z},\label{eq:ct-external-35}
\end{equation}
and, for a pair in separation class $k$,
\begin{equation}
P(I_{y,z}\cap I_{y',z})
 \le C(1+z)e^{-4L+2k}e^{-2z}
       e^{-c(k\wedge(L-k))^{1/4}}.\label{eq:ct-external-36}
\end{equation}
The convention for $L-d_*\le k<L$ is the one-center bound from that
lemma.  Fix $z_0$ in Rosen's admissible range.  After imposing the
source-safety event below we will choose a fixed $C_*$ and use the
truncated events
\begin{equation}
I_{y}^{\circ}:=I_{y,z_0}\cap
 \{\rho_LL+z_0\le\sqrt{2T_{y,L}^{1,m_0}}
                         \le\rho_LL+z_0+C_*\}\label{eq:ct-external-37}
\end{equation}
The notation in \eqref{eq:ct-external-37} is provisional until $C_*$ is chosen after
safety has been proved.  This order matters: because $X$ was extended
after $\tau$, the inequality
$T_{y,L}^{1,m_0}\le T_{y,L}^{\tau}$ can fail before safety is imposed.

\subsubsection{Keeping the source excursions before the stopping time}

The lower event was defined using a Brownian continuation after $\tau$.
We must therefore force the entire source forest to finish before $\tau$
without losing its probability on the scale $e^{-2L}$.  Let $S_y$ require every connector
$[B_{y,i},A_{y,i+1}]$, $i<m_0$, to hit $\Gamma_{y,1}$ before the
larger circle $\Gamma_{y,-1}$.  In the centered logarithmic cylinder,
the unrestricted connector endpoint subdensity and the safely killed
one, relative to Haar measure, are respectively
\[
q_\infty(u)={\sinh1\over\cosh1-\cos u},
\]
\begin{equation}
q_{\rm sf}(u)=\frac12+2\sum_{n\ge1}{\sinh n\over\sinh(2n)}\cos(nu)
 =\frac\pi2\sum_{k\in\mathbb Z}
       \operatorname {sech}\!\left({\pi\over2}(u+2\pi k)\right).\label{eq:ct-external-38a}
\end{equation}
Both are continuous and strictly positive, so
$\kappa_*:=\min_uq_{\rm sf}(u)/q_\infty(u)>0$.  Thus the safe
connector kernel dominates $\kappa_*$ times the unrestricted kernel
pointwise.  Iterating the strong Markov factorization, with the whole
later path functional left in the last factor, gives
\begin{equation}
P(I_{y,z}\cap S_y)\ge\kappa_*^{m_0-1}P(I_{y,z}).\label{eq:ct-external-38b}
\end{equation}
This is finite-measure domination against the entire later functional,
not multiplication of marginal probabilities.  Before endpoint
truncation put
\[
\widetilde Z_L^{\rm sf}:=\sum_{y\in F_L^0}{\bf1}_{I_{y,z_0}\cap S_y}.
\]
Equation \eqref{eq:ct-external-38b} preserves its first moment, whereas
$I_{y,z_0}\cap S_y\subset I_{y,z_0}$ preserves the pair upper bound
\eqref{eq:ct-external-36}.  Rosen's shell summation and Paley--Zygmund therefore give
\begin{equation}
E\widetilde Z_L^{\rm sf}\le C,\qquad
 P(\widetilde Z_L^{\rm sf}>0)\ge c_1>0.\label{eq:ct-external-38c}
\end{equation}
Rosen chooses $r_0$ so that $4h(r_{-1})<r_*$.  Since
$d(v,y)\le h_1/20$, there is a fixed $\chi>0$ such that
\begin{equation}
B_d(y,h(r_{-1}))\subset B_d(v,r_*-\chi).\label{eq:ct-external-38d}
\end{equation}
Every source block and every connector on $S_y$ lies in the left
side, hence ends before $\tau$.  This is the exact stopped-source
certificate used in the coupling below.  In particular, on $S_y$,
\begin{equation}
T_{y,L}^{1,m_0}\le T_{y,L}^{\tau}.\label{eq:ct-external-38e}
\end{equation}
Rosen's Theorem~1.6 gives the following upper tail on $F_L$; extending
its bound to bounded nonnegative $z$ only enlarges the constant:
\begin{equation}
P\left(\sup_{y\in F_L}\sqrt{2T_{y,L}^{\tau}}\ge\rho_LL+z\right)
       \le C(1+z)e^{-2z},\qquad0\le z\le\log L.\label{eq:ct-external-38f}
\end{equation}
Choose $C_*$ so large that the right side at $z=z_0+C_*$
is below $c_1/2$, and now define $I_y^\circ$ by \eqref{eq:ct-external-37}.  Equations
\eqref{eq:ct-external-38c}--\eqref{eq:ct-external-38f} give, with
\[
Z_L^{\rm sf}:=\sum_{y\in F_L^0}{\bf1}_{I_y^\circ\cap S_y},
\]
\begin{equation}
EZ_L^{\rm sf}\le C,\qquad P(Z_L^{\rm sf}>0)\ge c_1/2.\label{eq:ct-external-38}
\end{equation}
The constant $C_*$ and the safe count are now fixed.

\begin{lemma}[Continuum barrier inputs]\label{lem:external-barrier-input}
In the upper choice of radii, the first-crossing estimate
\eqref{eq:ct-external-28} holds.  For the complete centered traversal process,
the endpoint estimate \eqref{eq:ct-external-29} holds, including the truncated
barrier described in \eqref{eq:ct-external-31a}.  In the lower choice of radii,
on $I_y^\circ\cap S_y$ all first $m_0$ source excursions finish before
$\tau$, and the number of centers with this event satisfies
\eqref{eq:ct-external-38}.
\end{lemma}

\begin{proof}[Proof of Lemma~\ref{lem:external-barrier-input}]
The representative estimates are the packing and one-center bounds of
\cite[Proposition~2.5, (2.34)--(2.35)]{Rosen2022}, summed over the indicated
source shells.  The complete centered source calculation above proves
\eqref{eq:ct-external-29}, and \eqref{eq:ct-external-31a} checks the change of
barrier at a truncated horizon.  The lower construction combines
\cite[Lemmas~3.3--3.4]{Rosen2022} with the connector-kernel domination
\eqref{eq:ct-external-38b} and the terminal truncation.
All the first- and second-moment assertions \eqref{eq:ct-external-35}--\eqref{eq:ct-external-38} remain in Rosen's
original spherical geometry.  In particular, we do \emph{not} infer an
event-level Euclidean two-center theorem from Rosen's occupation-measure
argument in Section 8.  The next lemma instead maps the actual spherical
interfaces into the plane, so the spherical event and its pair law are
preserved exactly until the one-sided lattice coupling is made.
\end{proof}

\subsection{Comparison of smooth interface kernels}

Fix $\alpha\in(0,1)$.  Let
$(\Gamma_R^-,\Gamma_R^0,\Gamma_R^+)$ range over triples of
disjoint Jordan curves which, after translation, rotation, and scaling
by $R^{-1}$, together with their tubular-coordinate charts, belong to
a bounded $C^{4,\alpha}$ family.  Assume consecutive curves are
separated by distances comparable to $R$, and that a specified
conformal coordinate sends each triple to three concentric logarithmic
levels.  The interfaces used below are circles or uniformly
Möbius-transformed circles, so this hypothesis entails no loss.

For a boundary function put
\begin{equation}
\|f\|_{4,R}:=\sum_{j=0}^4R^j\|\partial_\tau^jf\|_\infty
             +R^{4+\alpha}[\partial_\tau^4f]_\alpha .\label{eq:ct-smooth-39a}
\end{equation}
A lattice starting state is a directed edge whose endpoints lie on opposite
sides of $\Gamma_R^0$; the lattice query starts from its post-crossing
endpoint and the Brownian query from the normal projection of that endpoint
onto $\Gamma_R^0$.  The side of a vertex is determined by membership in
the bounded component of the complement (vertices lying on the curve are
assigned to the exterior).  Let $Q_R^B$ be the Brownian selected-side
subprobability kernel, stopped at first contact with either competing curve,
and let $Q_R^S$ be the lattice selected-side subprobability kernel, stopped
at the first directed edge whose endpoint-side indicators differ for one of
the two competing curves.  A tangency, or any intersection of an interpolated
edge whose two endpoints remain on the same side, is not a lattice crossing.
Write $\pi_0,\pi_+$ for normal projection of the digital source
and selected target states.  For a function $f$ on the selected target
curve and a function $g$ on the middle curve, define the two differently
typed lifts by
\[
U_+f=f\circ\pi_+,
                 \qquad U_0g=g\circ\pi_0.
\]
\begin{lemma}[Comparison of interface kernels]
\label{lem:smooth-interface-discretization}
As an estimate from functions on the selected target curve to functions
on the digital source-edge space,
\begin{equation}
\|Q_R^SU_+f-U_0(Q_R^Bf)\|_\infty
       \le CR^{-1}\|f\|_{4,R}.\label{eq:ct-smooth-39b}
\end{equation}
If the Brownian side probability $p_R=Q_R^B1$ is constant on the
middle curve and lies in $[c,1-c]$, then, uniformly in the starting
angle,
\begin{equation}
Q_R^S1=p_R(1+O(R^{-1}))\label{eq:ct-smooth-39c}
\end{equation}
and the corresponding row-normalized angular kernels differ by
$O(R^{-1})$ on $C^4$ tests.

Consequently, let a chronological radial word consist of $N$
ordinary decisions between neighboring concentric interfaces at
scales $R_i\ge M$.  The next queried triple may depend on the already
exposed radial prefix, but not on an unexposed angular mark.  For a fixed
word, let $E_i$ be the directed-edge state space after its first $i$
letters and let $K^S_{i,w_{i+1}}(e,de')$ be the subkernel realizing the
next prescribed side.  If $p_i$ is the corresponding Brownian side
probability, then, for every current edge,
\begin{equation}
e^{-C/R_i}p_i\le K^S_{i,w_{i+1}}1(e)
          \le e^{C/R_i}p_i .\label{eq:ct-smooth-39cprime}
\end{equation}
After summing all compatible directed-edge refinements,
\begin{equation}
e^{-C\sum_iR_i^{-1}}P^B(w)
 \le P^S(w)\le
 e^{C\sum_iR_i^{-1}}P^B(w).\label{eq:ct-smooth-39d}
\end{equation}
The same inequality holds after integrating the initial edge against any
entrance law and after summing a disjoint family of complete words.  A
radial-prefix stopping rule is padded by cemetery letters of row mass one.
No assertion is made after conditioning on an edge list generated later by
the word.  Source-root and source-death letters are handled explicitly in
Lemma \ref{lem:buffered-transfer}.

\end{lemma}

\begin{proof}[Proof of Lemma~\ref{lem:smooth-interface-discretization}]
Let $\Omega_R$ be the annulus between the two competing curves and let
$u$ be the continuum Dirichlet solution with boundary data $f$ on the
selected component and zero on the other.  Uniform rescaled boundary
regularity gives
\begin{equation}
\|D^ju\|_{L^\infty(\overline\Omega_R)}
       \le CR^{-j}\|f\|_{4,R},\qquad0\le j\le4.\label{eq:ct-smooth-39e}
\end{equation}
The compact rescaled family has a tubular radius $\rho R$.  In a signed
normal chart $\Phi(s,t)$, with $t\ge0$ on $\Omega_R$, choose fixed
numbers $c_1,\ldots,c_5$ satisfying
\[
\sum_{a=1}^5c_a(-a)^q=1,
             \qquad q=0,1,2,3,4,
\]
and, for $-\rho R/5<t<0$, set
\[
\widetilde u(\Phi(s,t))=\sum_{a=1}^5c_a u(\Phi(s,-at)).
\]
Do this in the two disjoint collars and retain $\widetilde u=u$ on
$\overline\Omega_R$.  The five matching identities make
$\widetilde u$ a $C^4$ continuation through the boundary, with the
same scaled derivative bounds as \eqref{eq:ct-smooth-39e}.

Let $\Omega_R$ be understood with the same vertex-side convention and set
\[
T=\inf\{n\ge1:S_{n-1}\in\Omega_R,\ S_n\notin\Omega_R\}.
\]
Equivalently, $T$ is the first side-changing directed edge across one of
the two competing components; same-side contacts of the interpolated edge
are ignored.  For $n<T$, $S_n$ lies in the open annulus or on its inner
boundary (with the exterior convention above); it cannot lie on the outer
boundary.  At an open-annulus point
$\Delta\widetilde u(S_n)=\Delta u(S_n)=0$, while at an inner-boundary
lattice vertex the $C^4$ matching and continuity from the annulus give
\[
\Delta\widetilde u(S_n)
       =\lim_{\Omega_R\ni x\to S_n}\Delta u(x)=0.
\]
Thus Taylor expansion is legitimate on the full four-neighbor stencil and
its quadratic term vanishes in either case.  The linear and cubic terms
cancel by symmetry.
Thus
\begin{equation}
|(P-I)\widetilde u(S_n)|\le CR^{-4}\|f\|_{4,R},
             \qquad n<T.\label{eq:ct-smooth-39f}
\end{equation}
For large $R$, the two boundary components are more than two lattice
spacings apart and each has tubular radius $\rho R$.  Hence the exit edge
crosses exactly one component, its post-crossing endpoint is within one
lattice spacing of that component and has a unique normal projection.
Moreover $T$ is no larger than the exit time of a fixed $CR$-disk
containing the annulus.  Optional stopping of
$\lVert S_n-z_*\rVert^2-n$ gives $ET\le CR^2$, so the accumulated
defect in \eqref{eq:ct-smooth-39f} is $O(R^{-2})\|f\|_{4,R}$.  The terminal endpoint is
within one lattice unit of the uniquely crossed component; its value
under $\widetilde u$ differs from the projected boundary payoff by
$O(R^{-1})\|f\|_{4,R}$.  The post-crossing starting endpoint and its
projection onto $\Gamma_R^0$ have the same error.  Optional stopping of
$\widetilde u(S_{n\wedge T})$ minus its defect sum proves \eqref{eq:ct-smooth-39b}.
Taking $f=1$, using $p_R\in[c,1-c]$, and dividing by the row mass
proves \eqref{eq:ct-smooth-39c}.  Notice that the continuation is used only to make a
boundary-crossing stencil legal; harmonicity, or its boundary-continuity
consequence, is invoked only at a live vertex.

For completeness, the logarithmic annular probability behind \eqref{eq:ct-smooth-39c} is
also the following cited lattice estimate.  Hao--Li--Okada--Zheng,
\emph{Favorite sites for simple random walk in two and more dimensions},
\cite{HLOZ2025}, Lemma 2.1,
equation (2.5), states in full that, uniformly for
$r<\lVert x\rVert<R$,
\[
P_x(H_{\partial D(0,R)}<H_{\partial D(0,r)})
 ={\log(\lVert x\rVert/r)+O(r^{-1})\over\log(R/r)},
\]
and the complementary inner-hit probability is
\[
P_x(H_{\partial D(0,r)}<H_{\partial D(0,R)})
 ={\log(R/\lVert x\rVert)+O(r^{-1})\over\log(R/r)}.
\]
At fixed logarithmic gaps the main terms are bounded away from zero.

Finally fix a prescribed radial word and let $\nu_i^S$ be the finite
measure on its current directed-edge states after the first $i$ letters,
with every earlier edge integrated out.  The radial prefix fixes the next
interface and prescribed side before the current edge is integrated.
Brownian conformal radial symmetry makes its row mass $p_i$ independent
of the angular endpoint, while \eqref{eq:ct-smooth-39cprime} gives
\begin{equation}
e^{-C/R_i}p_i\|\nu_i^S\|
 \le\|\nu_i^SK^S_{i,w_{i+1}}\|
 \le e^{C/R_i}p_i\|\nu_i^S\|.\label{eq:ct-smooth-39iprime}
\end{equation}
Forward induction proves \eqref{eq:ct-smooth-39d}.  This remains valid for an adaptive
exploration because the radial prefix determines its next query, and for a
stopped exploration after adjoining the absorbing cemetery state.  It is a
forward marginal comparison, not a conditional law given future entrance
points, and it never applies \eqref{eq:ct-smooth-39b} to a rare future indicator.
\end{proof}

\subsection{Transfer of traversal events to the lattice}

Fix $A,C_0<\infty$ and a positive barrier intercept $z$ in the
admissible range above.  The upper first-crossing estimate below is uniform
for $0\le s\le C_0\log J$.  Choose $h=h_J$ so that
\begin{equation}
{h\over\log J}\longrightarrow\infty,\qquad h=o(J),\qquad
 M=e^h,\qquad n=J-h+O(1).\label{eq:ct-transfer-39}
\end{equation}
For the lower comparison let $v$ and $r_*$ be Rosen's starting
point and outer stopping radius.  Fix stereographic coordinates $\sigma$ with
$\sigma(v)=0$, put $R_*=2\tan(r_*/2)$, and dilate by
$A_J=e^J/R_*$:
\[
\Psi_J=A_J\sigma .
\]
Thus the spherical stopped path becomes, up to an increasing time
change, planar Brownian motion started at zero and stopped on
$\partial B(0,e^J)$.  For a spherical center $y$, let $\chi_y$ be
the rotated stereographic coordinate satisfying
\[
\chi_y(y)=0,\qquad |\chi_y(x)|=2\tan(d(x,y)/2).
\]
The lower-comparison physical interfaces are
\begin{equation}
\Gamma_{y,j}=\Psi_J\{x:|\chi_y(x)|=r_0e^{-j}\},
       \qquad-1\le j\le n.\label{eq:ct-transfer-39g}
\end{equation}
They are Euclidean circles, although their Euclidean centers need not be
the same.  In the conformal coordinate
$\chi_y\circ\Psi_J^{-1}$, they are exactly the concentric logarithmic
levels used by Rosen.  Their physical scale is comparable to
$e^{J-j}$; the $O(1)$ choice in $n=J-h+O(1)$ makes the deepest
scale comparable to $M$.

The linearly interpolated lattice path identifies which smooth curve a
lattice edge crosses, but every probabilistic cut is made at an integer
edge-index stopping time.  If $D_{y,j}$ is the bounded component inside
$\Gamma_{y,j}$, then, for adjacent levels and after an outer hit
$\eta_i$, set
\begin{equation}
\begin{aligned}
\sigma_i&=\inf\{n>\eta_i:S_{n-1}\notin D_{y,j},\ S_n\in D_{y,j}\},\\
\eta_{i+1}&=\inf\{n>\sigma_i:S_{n-1}\in D_{y,j-1},\
                                  S_n\notin D_{y,j-1}\}.
\end{aligned}\label{eq:ct-transfer-39h}
\end{equation}
A state records the directed crossing edge; interpolation is used only to
assign its projected $\chi_y$-angle.
Only portions belonging to the first $m_0$ source intervals
$[A_{y,i},B_{y,i}]$ in \eqref{eq:ct-external-32b} contribute traversal counts or terminal
bridges.  Inter-source connectors contribute only their guarded coarse
boundary decisions, never their interior occupation.  The whole segment
from $\sigma_i$ to $\eta_{i+1}$, including further contacts with the
inner interface, is one outward piece.  Recursing inside each parent
excursion gives the standard nested excursion tree.  The lower
\emph{radial word} $W_y=(W_{y,1},\ldots,W_{y,m_0})$ is a fixed collection
of $m_0$ centered words, one for each source interval, with all angles
erased; the complete guard record additionally contains the finitely many
connector boundary decisions.  If
$\ell_{y,i}\in\{1,-1\}$ records whether the connector after
$B_{y,i}$ first hits $\Gamma_{y,1}$ or $\Gamma_{y,-1}$, write
\begin{equation}
\overline W_y=(W_{y,1},\ell_{y,1},W_{y,2},\ldots,
                 \ell_{y,m_0-1},W_{y,m_0})\label{eq:ct-transfer-39hprime}
\end{equation}
for the complete conditioned guard word.  Its traversal coordinate
$T^S_{y,j}$ is the sum over these
words of completed crossings from $\Gamma_{y,j-1}$ to
$\Gamma_{y,j}$.  At the last
fixed number of levels, the interfaces are replaced through one
unit-width digital collar by concentric circles about the nearest lattice
point to $\Psi_J(y)$, of radii $M_y,eM_y$ with
$M_y\in[cM,CM]$.  We relabel $M_y$ as $M$; every later estimate is
uniform under this fixed-factor change and $\log M_y=h+O(1)$.
Define ${\cal R}_y$ by the chronological erasure map which replaces each
completed terminal segment, from its inward directed crossing of radius
$M$ through its next outward directed crossing of radius $eM$, by the
two endpoint edges.  It retains $\overline W_y$, all radial letters and
all crossing edges and projected angles outside the erased segments, but
it does not retain an erased segment's duration, interior crossings, path,
or occupation.  Inward connector pieces outside these erased segments are
exposed.  This quotient sigma-field is the intrinsic coarse sigma-field.

The upper application is instead defined directly on the lattice and is
entirely centered.  Write $N=e^J$, $\tau_N=\tau_{D(0,N)}$, and, for
each candidate cell center with $\|y\|<2N$, set
\begin{equation}
\sigma_y=\tau_{D(y,3N)}.\label{eq:ct-transfer-39k}
\end{equation}
The pathwise inclusions
\begin{equation}
D(0,N)\subset D(y,3N),\qquad \tau_N\le\sigma_y,
 \qquad L_{\tau_N}(x)\le L_{\sigma_y}(x)\label{eq:ct-transfer-39l}
\end{equation}
hold for every $x$.  For each one-cell probability we decompose the
extended path up to $\sigma_y$ using only circles centered at $y$,
and sum those probabilities only afterwards.  Replacing $N$ by $3N$
changes every logarithmic index by $O(1)$.  Use Rosen's formal horizon
$J$, its barrier $\rho_Jj+z+(j\wedge(J-j))^{1/4}$, and retain only
interfaces through $n$.  Let ${\cal U}_{J,s}$ be the union of these
centered lattice upper first-crossing events, with source shell at least
$s$, and let
${\cal G}^{S,J}_{y,n,z}$ be the corresponding no-crossing event through
$n-1$, and put
\[
B_n^J=\rho_Jn+z+(n\wedge(J-n))^{1/4}.
\]
On the detailed range $\|y\|>M$, its exact source index is
\begin{equation}
u_y=\log {3e^J\over\|y\|},\qquad
 k_y=\lceil u_y\rceil,\qquad m_y=k_y+1;\label{eq:ct-transfer-39m}
\end{equation}
centers with $\|y\|\le M$ belong to the separately treated near-start
range.  Thus $u_y\in(k_y-1,k_y]$ and $m_y-u_y\in[1,2)$.
For the lower application relabel the numerical depth $n$ as $L$,
use the genuine horizon-$L$ quantities $\rho_L,l_L$, and let
$I_y^{S,\circ}$ be the \emph{radial} lattice version of \eqref{eq:ct-external-37}: it retains
the traversal barrier and the bounded terminal band but drops every
angular screen.  Angular screens are used only in the spherical
Brownian certificate that supplies the second moment.

\begin{lemma}[Transfer of traversal events]\label{lem:buffered-transfer}
With the definitions above,
\begin{equation}
P({\cal U}_{J,s})\le J^{-2s/\log J+o(1)}+O(J^{-A}),\label{eq:ct-transfer-40}
\end{equation}
\begin{equation}
P\big({\cal G}^{S,J}_{y,n,z},
       \sqrt{2T^S_{y,n}}\in[a,a+1]\big)
 \le C(k_y+1)^C(2+B_n^J-a)n^{-2}e^{-a^2/(2n)},\label{eq:ct-transfer-41}
\end{equation}
uniformly for $k_y\le C_0\log J$ and every fixed central range
$\epsilon_0n\le a\le B_n^J+1$.  No such assertion is made in the
extinction range $a<\epsilon_0n$.  Most importantly, if
\[
Z_L^S=\sum_{y\in F_L^0}{\bf1}_{I_y^{S,\circ}},
\]
then
\begin{equation}
EZ_L^S\le C,\qquad P(Z_L^S>0)\ge c.\label{eq:ct-transfer-42}
\end{equation}
Only the Brownian proof of \eqref{eq:ct-transfer-42} uses the two-center estimate \eqref{eq:ct-external-36}.
No two-center lattice estimate is claimed or used.

\end{lemma}

\begin{proof}[Proof of Lemma~\ref{lem:buffered-transfer}]
There are two logically separate comparisons.  Rare upper and endpoint
events are compared directly, word by word, using Lemma
\ref{lem:smooth-interface-discretization}.  The global strong coupling is used
only in the safe direction ``strict Brownian lower root implies loose
lattice lower root.''  At no point is a Brownian bad-subkernel estimate
multiplied by an arbitrary future functional of the coupled lattice path.
In the geometric estimates below $L$ denotes the numerical depth $n$;
the superscript $J$ records when the upper barrier retains the larger
formal horizon $J$.

\smallskip\noindent
\underline{\textit{Step 1. Geometry and one global coupling.}}

The conformal image retains the original spherical excursion structure.  Stereographic
projection is conformal and sends spherical Brownian motion to planar
Brownian motion under a continuous increasing time change.  Moreover,
by the definition of $\chi_y$,
\[
\Psi_J\bigl(\partial B_d(y,h(r_0e^{-j}))\bigr)
       =\Gamma_{y,j}.
\]
Thus hitting orders, the complete nested excursion tree, all traversal
counts, and all $\chi_y$-angular marks are literally Rosen's spherical
objects.  In particular, \eqref{eq:ct-external-36} will be applied only before any lattice
comparison.

For the strong approximation we use Zaitsev, Estimates for the Strong
Approximation in Multidimensional Central Limit Theorem,
\cite{Zaitsev2003},
Theorem 2 and Corollary 1.  The form needed here is: for fixed
dimension $d$, centered independent vectors with identity covariance
whose laws lie in the analytic finite-exponential-moment class
${\cal A}_d(\tau)$ for a fixed $\tau\ge1$, one can couple them with independent standard
Gaussian vectors so that, for
$\Delta_N=\max_{k\le N}\lVert\sum_{i\le k}(X_i-Y_i)\rVert$,
\[
E\exp\{c\Delta_N/(\tau d^{7/2}\log^*d)\}
 \le \exp\{C d^{9/4+\alpha}\log^*(N/\tau^2)\},
\]
and hence
\begin{equation}
P(\Delta_N>C\log N+x)\le C e^{-cx},\qquad x\ge0,\label{eq:ct-transfer-30}
\end{equation}
for a fixed i.i.d. law.  The planar walk increment multiplied by $\sqrt2$
has identity covariance and bounded support, so the hypotheses apply.

Take the deterministic time cap
\[
T_J=e^{2J}J^{20}.
\]
Lemma \ref{lem:elementary-walk-estimates} and its Brownian analogue show that
the walk or Brownian motion remains inside the outer disk beyond $T_J$
with probability $e^{-cJ^{10}}$.  By \eqref{eq:ct-transfer-30}, after changing Brownian time
by the covariance factor $1/2$, there is a coupling for which
\begin{equation}
\sup_{t\le T_J}\lVert S^{\rm lin}_{2t}-B_t\rVert
       \le b_J:=CJ^8\label{eq:ct-transfer-31}
\end{equation}
except on an event of probability $e^{-cJ^8}$.  Brownian interpolation
between integer times has a smaller Gaussian error.

The crossing convention is an intrinsic deterministic function of the
finite lattice path.  For a curve with bounded component $D$, a lattice
crossing is recorded exactly when
$\mathbf1_{\{S_{n-1}\in D\}}\ne\mathbf1_{\{S_n\in D\}}$; a curve vertex
is assigned to the exterior, and same-side tangencies or chord contacts are
ignored.  The crossing angle is the normal projection of the post-crossing
endpoint.  At all used scales the endpoint lies in the uniform tubular
collar, so this projection is unique.  On the compact family of
centers under consideration, the maps $\chi_y\circ\Psi_J^{-1}$ and
their inverses have uniformly bounded scaled first two derivatives.
Consequently, if two physical paths are within $b$, then every
crossing of a nominal interface by one path is sandwiched between the
crossings by the other path of the two interfaces obtained by changing
its $\log|\chi_y|$-level by $Cb/R$, where $R$ is the physical
scale.  This is just the deterministic inclusion of the two
$b$-neighborhoods of a smooth separating curve; applying it
successively preserves chronological order.
The crossed directed edge retains the digital-boundary overshoot, which
costs at most two extra lattice units.  Thus
the only difference between the nominal Brownian word and the lattice
word comes from collars of relative logarithmic width
\begin{equation}
\delta_J=C\frac{b_J+1}{M}=O\left(\frac{J^8}{M}\right).\label{eq:ct-transfer-33}
\end{equation}

\smallskip\noindent
\underline{\textit{Step 2. Preserving the lower radial word.}}

The target is the pathwise inclusion \eqref{eq:ct-transfer-R4b}.
It suffices to exclude a set $D_y$ of collar failures whose conditional
probability is small for each complete capped radial word. Outside $D_y$,
collar closeness must give the same ordered traversal forest, and hence the
same counts $T_{y,j}^{1,m_0}$, for the two paths. The three guards below
prevent a wrong neighboring hit, a retreat before the intended hit, and a
reversal just after it, respectively.

Uniform path closeness alone does not preserve an excursion tree: a path
may enter a narrow collar, retreat, and thereby insert or delete a traversal.
We exclude precisely these retreats and wrong-side approaches.  Their total
conditional probability will be small for every prescribed radial word.  Work in the centered cylinder coordinate
\[
U_t=-\log\bigl(|\chi_y(X_t)|/r_0\bigr),
\]
after its conformal time change, so that $U$ is one-dimensional Brownian
motion and $\Gamma_{y,j}=\{U=j\}$.  For a stopping time $s$, write
$H_a[s]=\inf\{t\ge s:U_t=a\}$, with $\inf\varnothing=\infty$.  Put
\begin{equation}
\eta_J=8\delta_J,\qquad N_0=J^4,\label{eq:ct-transfer-R1}
\end{equation}
and take $J$ large enough that $0<\delta_J<\eta_J/4$ and
$\eta_J\le1/4$.

Here is the finite depth-$L$ word belonging to \eqref{eq:ct-external-32b}.  From
$S_{i,0}=A_{y,i}$, $K_{i,0}=1$, recursively set, when $K_{i,r}<L$,
\[
S_{i,r+1}=\inf\{t>S_{i,r}:U_t\in
             \{K_{i,r}-1,K_{i,r}+1\}\},\qquad
 K_{i,r+1}=U_{S_{i,r+1}}.
\]
When $K_{i,r}=L$, instead set
\begin{equation}
S_{i,r+1}=H_{L-1}[S_{i,r}],\qquad K_{i,r+1}=L-1;\label{eq:ct-transfer-R1a}
\end{equation}
this is a forced terminal return, and the motion below $\Gamma_{y,L}$
is deliberately not resolved.  Stop at
$N_i=\inf\{r\ge1:K_{i,r}=0\}$, so $B_{y,i}=S_{i,N_i}$.
For $i<m_0$ put
\[
C_i=\inf\{t>B_{y,i}:U_t\in\{-1,1\}\}.
\]
On $S_y$, $U_{C_i}=1$, $C_i=A_{y,i+1}$, and the connector is the
single prescribed sign $0\to1$.  Concatenating the $m_0$ block words
and these $m_0-1$ connector signs gives exactly $\overline W_y$ in
\eqref{eq:ct-transfer-39hprime}.  A block letter $j-1\to j$ pushes the next Ulam--Harris child,
a letter $j\to j-1$ pops it, and \eqref{eq:ct-transfer-R1a} closes a terminal leaf.  Hence
\begin{equation}
T_{y,1}^{1,m_0}=m_0,\qquad
 T_{y,j}^{1,m_0}=\#\{j-1\to j\hbox{ block letters}\},
                  \quad2\le j\le L.\label{eq:ct-transfer-R1b}
\end{equation}
All these times are recursively defined hitting times of closed sets;
the completed $m_0$-block word and the cap $|\overline W_y|\le N_0$
are therefore stopping-line measurable.

We next define the guards.  Consider first a competitive letter beginning
at a stopping time $s$, at level $k<L$, whose prescribed sign is
$\epsilon\in\{-1,1\}$.  Its end is
\[
t=\inf\{u>s:U_u\in\{k-1,k+1\}\},\qquad U_t=k+\epsilon.
\]
Define
\[
\begin{array}{lll}
 p=H_{k+\epsilon(1-\eta_J)}[s],
 &w=H_{k-\epsilon(1-\eta_J)}[s],
 &r=H_{k+\epsilon\eta_J}[p],\\
 e=H_{k+\epsilon(1+\eta_J)}[t],
 &o=H_{k+\epsilon\eta_J}[t].
\end{array}
\]
The bad events for this letter are
\begin{equation}
D^{\rm wrong}=\{w<t\},\qquad
 D^{\rm ret}=\{r<t\},\qquad
 D^{\rm post}=\{o<e\}.\label{eq:ct-transfer-R2}
\end{equation}
The last comparison is with the old-side collar
$k+\epsilon\eta_J$, not merely with the old nominal line $k$.
For a forced terminal return $L\to L-1$ beginning at $s$, set
\[
t=H_{L-1}[s],\quad p=H_{L-1+\eta_J}[s],\quad
 r=H_{L-\eta_J}[p],\quad
 e=H_{L-1-\eta_J}[t],\quad o=H_{L-\eta_J}[t],
\]
and use $D^{\rm ret}=\{r<t\}$ and
$D^{\rm post}=\{o<e\}$.  There is no wrong-side guard because the
path below level $L$ is not resolved.

The first source opening needs its own two guards.  The fixed annulus
defining $F_L^0$ gives a uniform initial margin.  Indeed,
$d(v,y)\le h_1/20$, convexity of $\tan$ on the relevant fixed small
interval, and $r_1=r_0e^{-1}$ imply
\[
|\chi_y(v)|=2\tan\!\left(\frac{d(v,y)}2\right)
 \le 2\tan\!\left(\frac{h_1}{40}\right)
 \le \frac{r_1}{20},
 \qquad U_0\ge1+\log20.
\]
(The lower annulus bound keeps $U_0$ finite.)  Thus, with
$c_{\rm src}:=\log20$, one has $U_0\ge1+c_{\rm src}$, and for
$\eta_J<c_{\rm src}$ the stopping time
$c_1=H_{1+\eta_J}[0]$ satisfies $c_1<A_{y,1}$.  Since the first
opening approaches level $1$ from the inner side, define
\begin{equation}
D^{A_1,{\rm pre}}=\{H_{2-\eta_J}[c_1]<A_{y,1}\},\qquad
 D^{A_1,{\rm post}}=
 \{H_{2-\eta_J}[A_{y,1}]<H_{1-\eta_J}[A_{y,1}]\}.\label{eq:ct-transfer-R2a}
\end{equation}
For $i\ge2$, the pre-opening guard at $A_{y,i}$ is exactly
$D^{\rm ret}$ for connector $i-1$, and its post-opening guard is
exactly $D^{\rm post}$ for that connector; they are named both ways
but counted once.  At every $B_{y,i}$, the last block symbol is the
competitive closing letter $1\to0$; its wrong, retreat, and post events
in \eqref{eq:ct-transfer-R2} are respectively the pre-closing and closing-recognition guards.
For $i<m_0$ the following symbol is the connector $0\to1$.  At the
last close $B_{y,m_0}$, retain its post-guard and expose it without any
following letter.  Let $D_y$ be the union of these events through the
last post-closing collar hit.  Every member of this union has the form
$\{H_a[S]<H_b[S]\}$ for a named stopping time $S$, so $D_y$ is
measurable at that finite stopping line.

Each guard is decided at a named stopping line. Strong Markov cancels the
likelihood of all later prescribed letters in the numerator and denominator
of its conditional ratio; for a post-guard the next prescribed sign is kept
in the calculation. We record these ratios for an admissible complete word
$w$.  Translation, reflection, gambler's ruin, and strong Markov at the
integer-line hits give, for every competitive letter,
\begin{equation}
\frac{P(D^{\rm wrong},\overline W_y=w)}{P(\overline W_y=w)}
       =\frac{\eta_J}{2-\eta_J},\qquad
 \frac{P(D^{\rm ret},\overline W_y=w)}{P(\overline W_y=w)}
       =\frac{\eta_J(1+\eta_J)}
          {(1-\eta_J)(2-\eta_J)}.\label{eq:ct-transfer-R2b}
\end{equation}
If the next symbol is competitive, then
\begin{equation}
\frac{P(D^{\rm post},\overline W_y=w)}{P(\overline W_y=w)}
 =\begin{cases}
    \eta_J^2,&\text{if the next sign continues the present direction},\\
    2\eta_J-\eta_J^2,&\text{if it reverses it}.
   \end{cases}\label{eq:ct-transfer-R2c}
\end{equation}
Indeed, after translating the new level to $1$,
$P_1(H_{\eta_J}<H_{1+\eta_J})=\eta_J$; from $\eta_J$, the
probabilities of the next exits at $2$ and $0$ are respectively
$\eta_J/2$ and $1-\eta_J/2$, while either prescribed next sign has
probability $1/2$.  If the successor is the forced return in \eqref{eq:ct-transfer-R1a}, or
there is no successor, the post ratio is simply $\eta_J$.  For a forced
return $L\to L-1$,
\begin{equation}
\frac{P(D^{\rm ret},\overline W_y=w)}{P(\overline W_y=w)}
       =\frac{\eta_J}{1-\eta_J}.\label{eq:ct-transfer-R2d}
\end{equation}
The first-opening pre ratio in \eqref{eq:ct-transfer-R2a} is also
$\eta_J/(1-\eta_J)$.  If $d_1=-1$ is its crossing direction and
$\epsilon_1$ is the first block sign, its post ratio is
$\eta_J^2$ when $\epsilon_1=d_1$, and
$2\eta_J-\eta_J^2$ otherwise.  The same table applies at every later
opening through its incoming connector.  In particular, at $B_{y,i}$,
$i<m_0$, the close has direction $-1$ and the next connector has sign
$+1$, so its post ratio is $2\eta_J-\eta_J^2$.  At the final close
$B_{y,m_0}$ it is exactly $\eta_J$, with no fictitious next sign.
No angular or future functional is inserted into these ratios. A union bound, without
any independence assertion among guards, gives
\begin{equation}
P(D_y\mid\overline W_y=w)
 \le8(|w|+2m_0+1)\eta_J
 \le\varepsilon_J,
 \qquad \varepsilon_J:=CJ^4\delta_J=o(J^{-A})\label{eq:ct-transfer-R3}
\end{equation}
for every $|w|\le N_0$ and every fixed $A$.  The estimate remains
valid after conditioning on barriers, traversal counts, or any other
function of $w$, but is not claimed after conditioning on Rosen's
angular screens.

For completeness, here is the deterministic implication used below.  Let
$\mathcal A_J$ be the finite set of all nominal integer levels and all
guard-collar levels occurring in \eqref{eq:ct-transfer-R2}--\eqref{eq:ct-transfer-R2a}, including the forced-return
collars.  Let $V$ be the corresponding log-radius of the linearly
interpolated lattice path, in the coupling time parametrization.  Its
nominal source blocks and truncated word are defined at integer edge indices
by the endpoint-side-change convention above: a letter is recorded only
when an edge's endpoints change sides, and interpolation along such an edge
is used only to locate its crossing and evaluate collar separation.
Same-side tangencies and chord contacts record no letter.  Assume the
following symmetric collar condition: whenever either $U_t$ or $V_t$ is
within $2\delta_J$ of $\mathcal A_J$, both log-radii are finite and
$|U_t-V_t|\le\delta_J$.  No comparison is required while both paths lie
strictly on the unresolved inner side of the deepest collars.  Suppose $S_y$,
$|\overline W_y|\le N_0$, and $D_y^c$ occur.  Before the target collar
of a competitive Brownian letter, $V$ cannot make a side-changing crossing
of the target nominal line, while $D^{{\rm wrong},c}$ prevents it from hitting the other
neighbor.  From the first target-collar hit until the nominal target,
$D^{{\rm ret},c}$ keeps $U$ beyond the old-side collar; hence an early
side-changing target crossing by $V$ cannot be followed by a reversal.  From the nominal
hit to the far collar, $D^{{\rm post},c}$ gives the same separation,
and at the far collar
\[
\epsilon\bigl(V-(k+\epsilon)\bigr)
       \ge\eta_J-\delta_J>0.
\]
Thus $V$'s next nominal neighboring line is exactly $k+\epsilon$,
and it has made no intervening push or pop.  The forced return \eqref{eq:ct-transfer-R1a} uses
the retreat and post guards only; deeper visits are ignored by both
truncated words.

At $A_{y,1}$, \eqref{eq:ct-transfer-R2a} first prevents a premature completed neighboring
letter and then forces $V$ across the source line.  The preceding
letter argument initializes every later $A_{y,i}$ through its connector.
Induction in chronological order on the active Ulam--Harris stack now
gives the same push/pop operation at every symbol.  At $B_{y,i}$ the
same stack empties; for $i<m_0$ the guarded connector has sign $0\to1$
and starts the next stack, while the unconditional final post-guard forces
$V$ across level $0$ at $B_{y,m_0}$.  Hence the two stopped ordered
forests are canonically isomorphic and
\begin{equation}
T_{y,j}^{1,m_0}(V)=T_{y,j}^{1,m_0}(U),\qquad1\le j\le L.\label{eq:ct-transfer-R4}
\end{equation}
Moreover the exact word length is
\begin{equation}
|\overline W_y|=2\sum_{j=2}^{L}T_{y,j}^{1,m_0}+2m_0-1
       \le CL^3<N_0,\label{eq:ct-transfer-R4a}
\end{equation}
on the lower barrier and terminal band.  On the single global coupling
event \eqref{eq:ct-transfer-31}, the symmetric collar condition holds.  Indeed, every
$2\delta_J$-neighborhood of a nominal or guard collar used by the
truncated word has physical radius at least $cM$.  If either path is in
such a neighborhood, physical $b_J$-closeness, the radial Lipschitz
inequality, chart distortion, and the $O(1)$ terminal-interface rounding
give
\[
|U_t-V_t|\le C(b_J+1)/M\le\delta_J
\]
after fixing the constant in \eqref{eq:ct-transfer-33}.  No log-radius comparison is asserted or
used while both paths are deeper than the terminal collars.  Source safety
keeps every collar test before global killing.  This condition suffices:
any premature target, wrong-neighbor, or reversal crossing occurs at a
member of $\mathcal A_J$, where collar closeness and
$\delta_J<\eta_J/4$ contradict the corresponding safe guard.  This also
covers a premature $L-1$ hit during a forced $L\to L-1$ return.
Consequently
\begin{equation}
I_y^{\circ}\cap S_y\cap D_y^c\cap\{\text{global coupling \eqref{eq:ct-transfer-31}}\}
       \subseteq I_y^{S,\circ}.\label{eq:ct-transfer-R4b}
\end{equation}
Only the radial consequence is transferred: Rosen's original angular
screens remain solely on the Brownian certificate.

\smallskip\noindent
\underline{\textit{Step 3. The upper first-crossing probability.}}

We now prove the rare upper estimates by direct lattice comparison.  Fix a target cell
and use the centered continuation \eqref{eq:ct-transfer-39k}--\eqref{eq:ct-transfer-39l}.  Put
\[
R_j=R_0e^{-j},\qquad R_0=3e^J,
\]
rounding every radius by at most two lattice units.  On the detailed-source
range $\|y\|>M$, define from the actual starting radius
\begin{equation}
u_y=\log {R_0\over\|y\|},\qquad
 k_y=\lceil u_y\rceil,\qquad m=k_y+1.\label{eq:ct-transfer-45a0}
\end{equation}
Then $u_y\in(k_y-1,k_y]$, $m-u_y\in[1,2)$, and the
deterministic start lies outside $D(y,R_{m-1})$.  This index differs by
only $O(1)$ from $\lfloor\log(e^J/\|y\|)\rfloor$, so it changes no
shell exponent.  Write
\begin{equation}
\begin{split}
 E_j^{\rm in}&=\{(a,b):a\notin D(y,R_j),\ b\in D(y,R_j),\ a\sim b\},\\
 E_j^{\rm out}&=\{(a,b):a\in D(y,R_j),\ b\notin D(y,R_j),\ a\sim b\}.
\end{split}\label{eq:ct-transfer-45a1}
\end{equation}
Every radial query retains its phase, level, Ulam--Harris stack and child
counters, population vector, and current directed crossing edge.

Potential-kernel optional stopping, uniformly in that edge, gives for
$k<l$
\begin{equation}
q^S_{k,l}\le C{k+1\over l},\qquad
 p_l^S(e)={l-1+O(R_l^{-1})\over l+O(R_l^{-1})},\qquad
 d_l^S(e)={1+O(R_l^{-1})\over l+O(R_l^{-1})}.\label{eq:ct-transfer-45a}
\end{equation}
Here $q^S_{k,l}$ includes the first completion at level $l$, while
$p_l^S$ and $d_l^S=1-p_l^S$ are the continuation and death rows
after return to level $l-1$.  Indeed, the two complete formulas from
Hao--Li--Okada--Zheng, \cite{HLOZ2025}, Lemma 2.1, equation (2.5), quoted in
Lemma \ref{lem:smooth-interface-discretization}, have respective logarithmic
gaps $l-1+O(R_l^{-1})$, $1+O(R_l^{-1})$, and
$l+O(R_l^{-1})$.  Thus both the continuation and the order-$1/l$
death probability are relative estimates.  Strong Markov at every
successive directed hit yields, for $N\ge1$,
\begin{equation}
P(T_{y,l}^{S,k\to0}\ge N\mid\hbox{initial directed edge})
 \le C{k+1\over l}\exp\left\{-{N\over l}+{CN\over lM}\right\}.\label{eq:ct-transfer-45b}
\end{equation}
In particular, for $X=\sqrt{2T_{y,m}}$ and $x\le Cm$,
\begin{equation}
P(X\in[x,x+1])
       \le Cm^C e^{-x^2/(2m)}.\label{eq:ct-transfer-45c}
\end{equation}
No independence of directed entrance edges is used in these scalar bounds.

A union bound over all fine centers would count many copies of the same
crossing.  We replace each center by a deterministic nearby representative
and enlarge its target disk so that the crossing count can only increase.  Put
\[
\phi(l)=\exp\{\tfrac14(l\wedge(J-l))^{1/4}\}.
\]
Fix constants
\[
0<c_{\rm sep}\le c_{\rm mesh}<C_\delta<\frac{e-1}{e+1}.
\]
For each $l$, choose a deterministic maximal
$c_{\rm sep}R_l/\phi(l)$-separated subset of $\mathbb Z^2$, with
$c_{\rm sep}$ chosen so that its covering radius (for lattice points) is
at most $c_{\rm mesh}R_l/\phi(l)$.  Given a candidate lattice center
$y$ in source shell $k$, choose a mesh point
$\widehat y=\widehat y(y,l)\in\mathbb Z^2$ with
\[
|y-\widehat y|\le c_{\rm mesh}R_l/\phi(l),\qquad
 \delta_l=C_\delta/\phi(l),\qquad R_j^\pm=(1\pm\delta_l)R_j.
\]
For all sufficiently large $J$, the gap
$(C_\delta-c_{\rm mesh})R_l/\phi(l)$ absorbs the at-most-two-unit
rounding of every radius on the used range.  Moreover
$\delta_l<(e-1)/(e+1)$, and hence
$R_{l-1}^->R_l^+$, so every buffered annulus below is nonempty.
Starting with the first entrance into
$D(\widehat y,R_k^+)$, and stopping on exiting
$D(\widehat y,R_0^+)$, let
$\widetilde T_{\widehat y,l}^{S,k\to0}$ count completed crossings from
$\partial D(\widehat y,R_{l-1}^-)$ to
$\partial D(\widehat y,R_l^+)$.  The inclusions
\begin{equation}
\begin{split}
 D(y,R_k)&\subset D(\widehat y,R_k^+),\\
 D(y,R_l)&\subset D(\widehat y,R_l^+),\\
 D(\widehat y,R_{l-1}^-)&\subset D(y,R_{l-1}),\\
 D(y,R_0)&\subset D(\widehat y,R_0^+)
\end{split}\label{eq:ct-transfer-45d0}
\end{equation}
show pathwise that
\begin{equation}
T_{y,l}^{S,k\to0}
       \le\widetilde T_{\widehat y,l}^{S,k\to0}.\label{eq:ct-transfer-45d1}
\end{equation}
Indeed every $y$-crossing begins outside the representative's buffered
outer circle, ends inside its buffered inner circle, and occurs after the
representative start and before its stop.  The same inclusions cover
$k=l-1$.  Area comparison for the mesh in a shell of diameter
$O(R_k)$ gives, for every $k\le l-1$,
\begin{equation}
\#\{\hbox{representatives at }(k,l)\}
 \le C e^{2(l-k)}\phi(l)^2.\label{eq:ct-transfer-45d}
\end{equation}
The buffered inclusions in Rosen (2.26)--(2.33) do not impose
$k\le\log J$.  For the four modified radii
$(R_k^+,R_{l-1}^-,R_l^+,R_0^+)$, the same potential-kernel
calculation gives, uniformly in all directed entrance edges,
\begin{equation}
\widetilde q_{k,l}^S\le C{k+1\over l},\qquad
 \widetilde p_l^S
 \le1-\frac{1-C/\phi(l)-C/M}{l}.\label{eq:ct-transfer-45d2}
\end{equation}
Consequently \eqref{eq:ct-transfer-45b} holds with $T$ replaced by
$\widetilde T_{\widehat y,l}^{S,k\to0}$, the additional exponent being
$CN/(l\phi(l))+CN/(lM)$.  Taking
$N=\alpha_{z,+}(l)^2/2$ there gives
\begin{equation}
P(\widetilde T_{\widehat y,l}^{S,k\to0}\ge N)
 \le C e^{-2l-(l\wedge(J-l))^{1/4}}e^{-2z}.\label{eq:ct-transfer-45e}
\end{equation}
On the used range $l\le J-(4\log J)^4$, the elementary exponent
calculation gives $l/\phi(l)=O(1)$ (the bounded early values are
absorbed into $C$), and $N/(lM)=o(1)$.  Multiplying \eqref{eq:ct-transfer-45d} and \eqref{eq:ct-transfer-45e}
and summing first in $l$, then in $k\ge s$, proves \eqref{eq:ct-transfer-40}, including
the range $\log J<k\le C_0\log J$.

\smallskip\noindent
\underline{\textit{Step 4. The centered endpoint probability.}}

The source population and its descendants must be compared together.
Conditioning on a later list of entrance edges would change the law of the
current radial decision.  We therefore compare the entire chronological
word, integrating each edge when it is first produced, and only then
condition under the ideal Galton--Watson law.  In addition to
the directed-edge/stack states in \eqref{eq:ct-transfer-45a1}, introduce a special deterministic
initial atom $\star_y$, representing $S_0=0$.  The first query is
\texttt{ROOT} (hit $D_m$ before leaving $D_0$), whose kernel maps
$\star_y$ to $E_m^{\rm in}$, or \texttt{FAILURE}, whose kernel maps
$\star_y$ to the cemetery state.  After every root, its descendants are
explored depth first: at an active node of level $j<n$, \texttt{BIRTH} is the
first hit of $D_{j+1}$ and pushes the next child, while \texttt{CLOSE} is the
first exit from $D_{j-1}$ and pops the node.  A level-$n$ leaf has a
forced return to level $n-1$, of row mass one.  After a root closes, the
current state lies in $E_{m-1}^{\rm out}$; the next source query is
another \texttt{ROOT}, mapping to $E_m^{\rm in}$, or the final \texttt{DEATH}, mapping
through $E_0^{\rm out}$ to the cemetery state.  Thus a word with
$t\ge1$ roots has the
physical order
\begin{equation}
\mathtt{ROOT},\operatorname{DFS}(1),\mathtt{ROOT},
 \operatorname{DFS}(2),\ldots,\mathtt{ROOT},
 \operatorname{DFS}(t),\mathtt{DEATH}.\label{eq:ct-transfer-45f0}
\end{equation}
The word stops at the first barrier violation, endpoint decision, death,
or a fixed deterministic cap $Q_*=C_*J^3$, chosen above
$3+2\sum_{j=m}^n(T_{y,j}+1)$ on the coefficient-two barrier,
and is then padded by cemetery letters.  It contains every terminating
close and the final source death, so there is no residual future condition.

Let $P_*$ be the angle-free ideal law.  Its first root and failure
probabilities are $u_y/m$ and $(m-u_y)/m$; subsequent root and death
probabilities are $(m-1)/m$ and $1/m$; every descendant
birth/close decision has probability $1/2$.  Potential-kernel optional
stopping, with every output edge summed but every current directed edge
retained, gives the four lattice source rows
\begin{equation}
\begin{aligned}
 r_{0,y}^S
   &=\frac{u_y+O(R_m^{-1})}{m+O(R_m^{-1})},
 & f_{0,y}^S
   &=\frac{m-u_y+O(R_m^{-1})}{m+O(R_m^{-1})},\\
 r_m^S(e)
   &=\frac{m-1+O(R_m^{-1})}{m+O(R_m^{-1})},
 & d_m^S(e)
   &=\frac{1+O(R_m^{-1})}{m+O(R_m^{-1})}.
\end{aligned}\label{eq:ct-transfer-45f-src}
\end{equation}
Here the first pair starts from $\star_y$, while the second pair is
uniform in $e\in E_{m-1}^{\rm out}$.  Since every candidate satisfies
$\|y\|<2e^J$,
\[
u_y\ge\log(3/2),\qquad m-u_y\in[1,2),\qquad m-1\ge1.
\]
Also $R_m=e^{u_y-m}\|y\|\ge e^{-2}\|y\|\ge e^{-2}M$ on the detailed
range.  Consequently each of the four displayed rows, including the two
order-$1/m$ FAILURE and DEATH rows, equals its ideal row times
$1+O(M^{-1})$.  Ordinary birth/close rows are
$(1/2)(1+O(R_{j+1}^{-1}))$.

Fix a complete radial word $w$, but no directed-edge refinement.  If
$\nu_r$ is the edge-valued finite measure after its first $r$ letters,
the next ideal row $p_r$ is determined by the radial prefix, and the
uniform row estimates give
\[
e^{-C/M}p_r\|\nu_r\|\le\|\nu_{r+1}\|
       \le e^{C/M}p_r\|\nu_r\|.
\]
Forward induction, with all compatible entrance and exit edges integrated,
therefore gives
\begin{equation}
e^{-C|w|/M}P_*(w)\le P_S(w)
       \le e^{C|w|/M}P_*(w).\label{eq:ct-transfer-45f}
\end{equation}
This is not a conditional comparison given a future edge list.

Under $P_*$, the source population has the exact law
\begin{equation}
\pi_{m,u_y}(0)=1-{u_y\over m},\qquad
 \pi_{m,u_y}(t)={u_y\over m}
       \left({m-1\over m}\right)^{t-1}{1\over m},
       \quad t\ge1,\label{eq:ct-transfer-45g0}
\end{equation}
and, conditional on $t$, the $t$ descendant trees are independent
critical Galton--Watson trees with
$P(\xi=r)=2^{-r-1}$.  Hence, for the disjoint union
$\mathscr U(q)$ of complete barrier-respecting words ending in
$\sqrt{2T_n}\in[q,q+1]$,
\begin{equation}
P_S(\mathscr U(q))
 \le e^{CJ^3/M}\sum_{t\ge1}\pi_{m,u_y}(t)
       P_t^{\rm GW}(\mathscr A_t^{\rm GW}(q)).\label{eq:ct-transfer-45g}
\end{equation}
Only the ideal law on the right has been disintegrated.  Moreover, for
$x\le Cm$,
\begin{equation}
\sum_{\sqrt{2t}\in[x,x+1]}\pi_{m,u_y}(t)
 \le C(1+x)m^{-1}e^{-x^2/(2m)+O(1)}
 \le Ce^{-x^2/(2m)}.\label{eq:ct-transfer-45gprime}
\end{equation}
Apply Rosen Appendix Theorem 9.1(a), quoted in \eqref{eq:ct-external-30}, to the
ideal Galton--Watson factor in \eqref{eq:ct-transfer-45g}, group $\sqrt{2t}$ into unit bands,
and use \eqref{eq:ct-external-31}--\eqref{eq:ct-external-31a}.  This proves \eqref{eq:ct-transfer-41}.  The $t=1$ term is covered by
the extra-ancestor argument following \eqref{eq:ct-external-30}, while $t=0$ cannot reach a
positive central endpoint.  Since $J^3/M=o(J^{-A})$, the comparison is
relative even for the rare endpoint event.

\smallskip\noindent
\underline{\textit{Step 5. The lower count and the terminal interfaces.}}

We apply the Brownian second moment before transferring any event to the lattice.  We retain Rosen's original event
$I_y^\circ$ and its original angular screens exactly as written in
\eqref{eq:ct-external-33}--\eqref{eq:ct-external-37}, and intersect it with the source-safety event $S_y$.
No tightened or halved screen is used.  Put
$I_y^{\rm sf}=I_y^\circ\cap S_y$.  By \eqref{eq:ct-external-38},
\begin{equation}
P\left(\sum_y1_{I_y^{\rm sf}}>0\right)\ge c_* .\label{eq:ct-transfer-46}
\end{equation}
The screens occur only in this Brownian second-moment certificate; the
lattice target discards all angular screens.

Let $\widehat I_y^B$ be the Brownian event obtained from
$I_y^\circ$ by deleting every angular screen and retaining only the
radial barrier, the bounded endpoint band, and the fixed $m_0$-source
definition.  Rosen Appendix Theorem 9.1(a), followed by the same source
and barrier calculation as \eqref{eq:ct-external-29}, gives
\begin{equation}
P(\widehat I_y^B)\le Ce^{-2L},
             \qquad \sum_yP(\widehat I_y^B)\le C.\label{eq:ct-transfer-47}
\end{equation}
Since
$I_y^{\rm sf}\cap D_y\subseteq
  \widehat I_y^B\cap S_y\cap D_y$, estimate \eqref{eq:ct-transfer-R3} may be summed over
complete safe radial words before any angular screen is imposed.  Hence
\begin{equation}
\sum_yP(I_y^{\rm sf}\cap D_y)
 \le\varepsilon_J\sum_yP(\widehat I_y^B\cap S_y)
 \le\varepsilon_J\sum_yP(\widehat I_y^B)=o(1).\label{eq:ct-transfer-47a}
\end{equation}
Let ${\cal K}_J$ be the single global coupling/time-cap event \eqref{eq:ct-transfer-31}.
The pathwise implication \eqref{eq:ct-transfer-R4b} gives
\[
\{Z_L^{\rm sf}>0\}\cap{\cal K}_J
       \setminus\bigcup_y(I_y^{\rm sf}\cap D_y)
       \subseteq\{Z_L^S>0\}.
\]
Consequently, paying the coupling failure only once,
\begin{equation}
P(Z_L^S>0)\ge P(Z_L^{\rm sf}>0)
   -\sum_yP(I_y^{\rm sf}\cap D_y)-P({\cal K}_J^c)
 \ge c_*-o(1).\label{eq:ct-transfer-47b}
\end{equation}
Source safety and \eqref{eq:ct-external-38d} keep every block, connector, and final exposed
post-closing collar a fixed positive distance inside the stopping boundary,
so the close lattice path is not killed before the transferred blocks
finish.  No Brownian angular screen is transferred.
This is the only use of the two-center estimate \eqref{eq:ct-external-36}, and it occurs
entirely on the spherical Brownian side.

For the lattice first moment, extend the path if necessary and decompose
its first $m_0$ source blocks, treating every connector as an integrated
subkernel of row mass at most one.  For a fixed collection of block words,
apply the edge-marginal induction \eqref{eq:ct-smooth-39iprime} to all competitive letters,
uniformly over the entrance law supplied by each preceding connector.
Forced terminal returns have row mass one.  Thus, if $q$ is the number
of competitive decisions,
\[
P_S(\hbox{these block words})\le e^{CJ^3/M}2^{-q}.
\]
Summing the disjoint words obeying the lower barrier and terminal band gives
the critical geometric Galton--Watson probability with exactly $m_0$
ancestors.  Rosen Appendix Theorem 9.1(a), followed by the fixed-source
calculation already used in \eqref{eq:ct-transfer-47}, gives
\begin{equation}
P(I_y^{S,\circ})\le Ce^{-2L},\qquad EZ_L^S\le C.\label{eq:ct-transfer-47cprime}
\end{equation}
This proves \eqref{eq:ct-transfer-42}; in particular, the lower first moment does not use the
upper source-population estimate \eqref{eq:ct-transfer-45c}.

Finally we justify the terminal switch used in the definition of
${\cal R}_y$.  On the compact center family, with
$F_y=\sigma\circ\chi_y^{-1}$,
\[
F_y(w)=F_y(0)+a_yw+b_yw^2+O(|w|^3),
       \qquad0<c\le|a_y|\le C.
\]
At $r_L\asymp e^{-J+h}$, dilation by $A_J\asymp e^J$ shows that
the actual deepest image circle is at Hausdorff distance
\begin{equation}
O(A_Jr_L^2)=O(e^{-J+2h})=o(1)\label{eq:ct-transfer-47c}
\end{equation}
from its affine concentric circle.  Rounding the center and the digital
boundary costs $O(1)$, or relative width $O(M^{-1})$, and is one
additional collar in the lower guard construction \eqref{eq:ct-transfer-R1}--\eqref{eq:ct-transfer-R4}.  The induced
angular conjugacy, after removing the fixed rotation
$a_y/|a_y|$, is $C^4$-close to the identity by
$O(e^{-J+h})$; its Haar-density defect is therefore absorbed by
$M^{-7/24}\log M+M^{-1}$.  Thus the outer $\chi_y$-rotation kernels
and the terminal Euclidean rotation kernels use the same common center
up to the stated error.  All lattice
events were defined before introducing the auxiliary coupling and are
measurable in ${\cal R}_y$.  We may therefore discard the Brownian
path after \eqref{eq:ct-transfer-47b}; every later conditional law is the original
simple-random-walk law given ${\cal R}_y$.
\end{proof}
